\documentclass[a4paper, serif, 10pt]{moderncv}

%% moderncv
% style and color
\moderncvstyle{banking}
\moderncvcolor{black}

%% page layout/margins
\usepackage[top=2.5cm, bottom=2.5cm, left=1.5cm, right=1.5cm]{geometry}

\nopagenumbers{}

%% Babel config for Brazilian Portuguese language
% ref:
% - https://latex3.github.io/babel/guides/locale-portuguese.html
\usepackage[brazilian]{babel}

%% fontspec
\usepackage{fontspec}

\IfFontExistsTF{Linux Libertine}{
  \setmainfont[Ligatures={TeX, Common}, Numbers=OldStyle]{Linux Libertine}
}{}
\IfFontExistsTF{Linux Libertine O}{
  \setmainfont[Ligatures={TeX, Common}, Numbers=OldStyle]{Linux Libertine O}
}{}

%% CTeX
% CJK support, used to show CJK characters in the resume
%
% - fontset=none: disable builtin fontset but instead set the CJK font manually
% - heading=false: disable ctex heading
% - punct=kaiming: use kaiming punctuations styles for CJK
% - scheme=plain: use plain scheme, do not override \normalsize font size
% - space=auto: space settings for CJK characters
%
% ref:
% - http://ctan.mirrorcatalogs.com/language/chinese/ctex/ctex.pdf
\usepackage[UTF8, heading=false, punct=kaiming, scheme=plain, space=auto]{ctex}

\IfFontExistsTF{Noto Serif CJK SC}{
  \setCJKmainfont{Noto Serif CJK SC}
}{}
\IfFontExistsTF{Noto Sans CJK SC}{
  \setCJKsansfont{Noto Sans CJK SC}
}{}

%% line spacing
\usepackage{setspace}
\setstretch{1.125}

%% URL styling - use normal text instead of monospace
\urlstyle{same}

% Auto-underline all links
% Defer until \begin{document} so hyperref is already loaded
\AtBeginDocument{%
  \let\oldhref\href
  \renewcommand{\href}[2]{\oldhref{#1}{\underline{#2}}}%
}

%% Patch \httplink and \httpslink to handle full URLs
% The original moderncv macros always prepend http:// or https:// to URLs.
% This causes issues when the URL already has a protocol (e.g., https://example.com)
% resulting in malformed URLs like https://https://example.com.
% This patch checks if the URL already contains "://" and uses it directly if so.
% Note: We use \str_set:Nx (x-expansion) to expand macros like \@homepage first.
\ExplSyntaxOn
\str_new:N \g_yamlresume_url_str

\RenewDocumentCommand{\httpslink}{O{}m}{
  \str_set:Nx \g_yamlresume_url_str {#2}
  \str_if_in:NnTF \g_yamlresume_url_str {://}
    { \url{#2} }
    { \url{https://#2} }
}

\RenewDocumentCommand{\httplink}{O{}m}{
  \str_set:Nx \g_yamlresume_url_str {#2}
  \str_if_in:NnTF \g_yamlresume_url_str {://}
    { \url{#2} }
    { \url{http://#2} }
}
\ExplSyntaxOff

\name{Lucas Almeida}{}
\title{Engenheiro de Software Sênior | Full Stack \& Cloud}
\phone[mobile]{+55 11 98765-4321}
\email{lucas.almeida.dev@email.com}
\homepage{https://lucasalmeida.dev}

\address{Avenida Paulista, 1000, São Paulo, SP, Brasil, 01310-100}{}{}

\extrainfo{{\small \faGithub}\ \href{https://github.com/lucasalmeida-dev}{@lucasalmeida-dev} {} {} {} • {} {} {}
{\small \faLinkedin}\ \href{https://www.linkedin.com/in/lucasalmeida-dev}{@lucasalmeida-dev}}

\begin{document}

\maketitle

\section{Informações Básicas}

\cvline{}{Engenheiro de software com mais de 5 anos de experiência no desenvolvimento de aplicações web escaláveis e na construção de arquiteturas de microsserviços.
%
\begin{itemize}
\item Especialista em Node.js, React e Python, com atuação em ambientes cloud como AWS, Kubernetes e Docker
\item Experiência em liderança técnica, mentoria de desenvolvedores e colaboração com times multidisciplinares
\item Foco em qualidade de código, testes automatizados, CI/CD e boas práticas de segurança
\item Buscando novos desafios em projetos de alto impacto e inovação tecnológica
\end{itemize}}
\section{Formação}

\cventry{fev. de 2014–dez. de 2018}
        {Bacharelado, Ciência da Computação, Nota: 8.7}
        {Universidade de São Paulo}
        {\url{https://www.usp.br/}}
        {}
        {\begin{itemize}
\item Desenvolvimento de uma base sólida em computação, matemática e desenvolvimento de software
\item Participação em projetos de pesquisa na área de aprendizado de máquina, com bolsa de iniciação científica
\item Atuação como monitor da disciplina de Programação Orientada a Objetos por dois semestres
\end{itemize}
\textbf{Cursos}: Algoritmos e Estruturas de Dados, Programação Orientada a Objetos, Engenharia de Software, Sistemas Operacionais, Redes de Computadores, Banco de Dados}
\section{Experiência Profissional}

\cventry{jan. de 2023–Atual}
        {Engenheiro de Software Sênior}
        {CloudNova Tecnologia}
        {\url{https://www.cloudnova.com.br}}
        {}
        {\begin{itemize}
\item Liderou o redesenho da plataforma principal com arquitetura de microsserviços, resultando em 40\% de melhoria na escalabilidade
\item Implementou pipelines de CI/CD com GitHub Actions e ArgoCD, reduzindo o tempo de deploy de 45 para 10 minutos
\item Orientou 5 desenvolvedores júnior e plenos, conduzindo code reviews e definindo padrões de qualidade
\item Colaborou com as equipes de produto e design para priorizar e entregar funcionalidades de alto valor
\end{itemize}
\textbf{Palavras-chave}: Node.js, TypeScript, Kubernetes, AWS, Microsserviços, CI/CD}

\cventry{mar. de 2020–dez. de 2022}
        {Engenheiro de Software}
        {Inova Tech Soluções}
        {\url{https://www.inovatech.com.br}}
        {}
        {\begin{itemize}
\item Desenvolveu e manteve aplicações web com React, Node.js e PostgreSQL, atendendo mais de 200 mil usuários ativos
\item Implementou testes automatizados com Jest e Cypress, elevando a cobertura de testes para 85\%
\item Refatorou módulos legados, melhorando a performance e reduzindo o consumo de recursos em 30\%
\item Atuou em squads ágeis, participando de cerimônias, planejamentos e retrospectivas
\end{itemize}
\textbf{Palavras-chave}: React, Node.js, PostgreSQL, Docker, Testes Automatizados, Agile}

\cventry{fev. de 2019–fev. de 2020}
        {Desenvolvedor Full Stack Júnior}
        {AzulCode Startup}
        {\url{https://www.azulcode.com.br}}
        {}
        {\begin{itemize}
\item Auxiliou no desenvolvimento de um MVP para um aplicativo de mobilidade urbana usando React Native e Express
\item Criou e consumiu APIs REST, implementando autenticação e integração com serviços de terceiros
\item Participou de sessões de pair programming e contribuiu para a melhoria contínua do código
\end{itemize}
\textbf{Palavras-chave}: React Native, Express, JavaScript, REST APIs}
\section{Idiomas}

\cvline{Português}{Nativo ou Bilíngue \hfill \textbf{Palavras-chave}: Português brasileiro}
\cvline{Inglês}{Proficiência Profissional Fluente \hfill \textbf{Palavras-chave}: TOEFL iBT 98}
\cvline{Espanhol}{Proficiência Profissional Limitada \hfill \textbf{Palavras-chave}: Comunicação básica}
\section{Competências}

\cvline{Linguagens de Programação}{Especialista \hfill \textbf{Palavras-chave}: JavaScript, TypeScript, Python, Go, SQL}
\cvline{Desenvolvimento Web}{Especialista \hfill \textbf{Palavras-chave}: React, Node.js, Express, HTML5, CSS3}
\cvline{Cloud \& DevOps}{Avançado \hfill \textbf{Palavras-chave}: AWS, Docker, Kubernetes, Terraform, GitHub Actions}
\cvline{Banco de Dados}{Avançado \hfill \textbf{Palavras-chave}: PostgreSQL, MySQL, MongoDB, Redis}
\cvline{Arquitetura de Software}{Avançado \hfill \textbf{Palavras-chave}: Microsserviços, REST APIs, GraphQL, Clean Architecture}
\section{Prêmios}

\cventry{dez. de 2022}
        {Inova Tech Soluções}
        {Prêmio Inovação em Engenharia de Software}
        {}
        {}
        {Reconhecimento concedido pela automação do processo de testes e deploy, que reduziu drasticamente o tempo de entrega da squad.}

\cventry{ago. de 2017}
        {Conselho Nacional de Desenvolvimento Científico e Tecnológico (CNPq)}
        {Bolsa de Iniciação Científica}
        {}
        {}
        {Bolsa de pesquisa para o projeto "Aplicação de redes neurais na classificação de imagens médicas", desenvolvido na Universidade de São Paulo.}
\section{Certificados}

\cventry{mar. de 2024}
        {Amazon Web Services}
        {AWS Certified Solutions Architect - Associate}
        {\url{https://aws.amazon.com/certification/}}
        {}
        {}

\cventry{set. de 2022}
        {Cloud Native Computing Foundation}
        {Certified Kubernetes Application Developer (CKAD)}
        {\url{https://www.cncf.io/certification/ckad/}}
        {}
        {}
\section{Projetos}

\cventry{abr. de 2021–out. de 2021}
        {Plataforma web para facilitar a adoção responsável de animais de estimação, conectando ONGs a tutores em potencial.}
        {Petinder - Adoção de Animais}
        {\url{https://github.com/lucasalmeida-dev/petinder}}
        {}
        {\begin{itemize}
\item Desenvolvido com React, Node.js e PostgreSQL, incluindo painel administrativo e sistema de busca com filtros
\item Implementou autenticação JWT, envio de e-mails transacionais e integração com mapas
\item Projeto de código aberto com mais de 300 estrelas no GitHub
\end{itemize}
\textbf{Palavras-chave}: React, Node.js, PostgreSQL, Open Source}

\cventry{jan. de 2022–dez. de 2022}
        {Ferramenta interna para gestão de feedback entre colaboradores, com metas e reconhecimento.}
        {Plataforma de Feedback Contínuo}
        {\url{https://github.com/lucasalmeida-dev/feedback360}}
        {}
        {\begin{itemize}
\item Criado com Next.js, Prisma e GraphQL para atender cerca de 3 mil colaboradores
\item Permitiu a criação de ciclos de avaliação personalizáveis e relatórios gerenciais automatizados
\item Reduziu em 50\% o tempo gasto com processos manuais de avaliação de desempenho
\end{itemize}
\textbf{Palavras-chave}: Next.js, GraphQL, Prisma, TypeScript}
\section{Interesses}

\cvline{Tecnologia}{Inteligência Artificial, Computação em Nuvem, Open Source}
\cvline{Esportes}{Futebol, Corrida, Musculação}
\cvline{Música}{Violão, MPB, Jazz}
\section{Voluntariado}

\cventry{jan. de 2019–dez. de 2022}
        {Instrutor Voluntário de Programação}
        {Código do Bem}
        {\url{https://www.codigodobem.org.br}}
        {}
        {\begin{itemize}
\item Ministrou cursos introdutórios de lógica de programação e desenvolvimento web para jovens de baixa renda
\item Auxiliou na elaboração de material didático e na mentoria de projetos finais
\item Participou de mutirões de ensino com impacto direto em mais de 200 estudantes
\end{itemize}}

\end{document}